%% kapa: exam=1 round=24 year=2013 period=1 form=A subject=economics subject_ko=경제원론 seq=26 number=66
%% source: Q-Net 감정평가사 기출문제 (cst00309) artlSeq=5073802 file=extracted/1008/5073802/2013년 제24회 감정평가사 1차시험 기출문제/1교시(A형).hwp
%% conversion: kordoc 4.24.51 -> tools/gen1.py (auto); review=proofread-sonnet; flags=none
\begin{question}{66}
모든 시장이 완전경쟁적인 B국의 총생산량이 $Y=AL^{ 1- \theta } K^{ \theta }$으로 결정될 때, 다음 설명 중 옳지 않은 것은? (단, $Y$ 는 총생산량, $A$ 는 총요소생산성, $L$ 은 노동투입량, $K$ 는 자본투입량, 0<$\theta$<1이다.)
\begin{choices}
  \item 노동소득분배율은 $1- \theta$이다.
  \item 자본소득분배율은 $\theta$이다.
  \item 총생산함수는 규모에 대한 수익불변이다.
  \item 총요소생산성이 증가하면, 노동소득분배율 대비 자본소득분배율의 상대적 비율이 증가한다.
  \item 생산요소가 증가하지 않을 경우, 총생산량의 증가율은 총요소생산성 증가율에 의해서 결정된다.
\end{choices}
\end{question}
\answer{4}
